\chapter{Wegweiser zu den Hauptresultaten}
\label{b47guide:ch:roadmap}

\section{Welche Voraussetzung liefert welches Resultat?}

Die Verweise führen zu den vollständigen Aussagen mit ihren Voraussetzungen.
Alle Teilmengen tragen die jeweilige Teilraummetrik.

\begin{center}
\small
\begin{tabularx}{\linewidth}{@{}p{.30\linewidth}X@{}}
\toprule
Ausgangslage & Resultat und Fundstelle\\
\midrule
Punkt im Abschluss & Genau die Grenzwerte von Folgen aus der Menge;
Satz~\ref{b47guide:thm:sequential-closure}.\\[4pt]
Cauchy-Folge mit konvergenter Teilfolge & Die ganze Folge konvergiert;
Satz~\ref{b47guide:thm:cauchy-subsequence}.\\[4pt]
Abgeschlossener Teilraum eines vollständigen Raums & Vollständigkeit;
Satz~\ref{b47guide:thm:closed-complete}.\\[4pt]
Geschachtelte abgeschlossene Mengen mit Durchmessern gegen null &
Cantorsches Vollständigkeitskriterium;
Satz~\ref{b47guide:thm:cantor}.\\[4pt]
Totale Beschränktheit & Jede Folge besitzt eine Cauchy-Teilfolge;
Satz~\ref{b47guide:thm:tb-cauchy}.\\[4pt]
Vollständigkeit und totale Beschränktheit & Kompaktheit und
Folgenkompaktheit; Satz~\ref{b47guide:thm:compact-equivalence}.\\[4pt]
Vollständiger Umraum & Relative Kompaktheit ist totale Beschränktheit;
Satz~\ref{b47guide:thm:relative-compact}.\\[4pt]
Stetigkeit auf kompaktem Raum & Gleichmäßige Stetigkeit;
Satz~\ref{b47guide:thm:heine-cantor}.\\[4pt]
Dichter Definitionsraum und vollständiger Zielraum &
Gleichmäßig stetige Fortsetzung;
Satz~\ref{b47guide:thm:uniform-extension}.\\[4pt]
Beliebiger metrischer Raum & Eindeutig bestimmte Vervollständigung
bei festgelegter Einbettung;
Satz~\ref{b47guide:thm:completion-existence}.\\[4pt]
Kontraktion auf nichtleerem vollständigem Raum & Eindeutiger Fixpunkt,
konvergente Iteration und Fehlerkontrolle;
Sätze~\ref{b47guide:thm:banach} und~\ref{b47guide:thm:banach-errors}.\\[4pt]
Vollständigkeit & Bairescher Kategoriensatz;
Satz~\ref{b47guide:thm:baire}.\\[4pt]
Gleichgradige Stetigkeit und punktweise relative Kompaktheit &
Arzelà--Ascoli bei kompaktem Definitionsraum und vollständigem Zielraum;
Satz~\ref{b47guide:thm:ascoli}.\\
\bottomrule
\end{tabularx}
\end{center}

\section{Weitere metrische Konstruktionen und Kriterien}

\begin{center}
\small
\begin{tabularx}{\linewidth}{@{}>{\raggedright\arraybackslash}p{.36\linewidth}>{\raggedright\arraybackslash}X@{}}
\toprule
Ausgangslage & Resultat und Fundstelle\\
\midrule
Reellwertige Lipschitz-Abbildung auf einer beliebigen Teilmenge &
Fortsetzung mit derselben Lipschitz-Konstante;
Satz~\ref{b47guide:thm:mcshane-whitney}.\\[4pt]
Gleichmäßige Cauchy-Bedingung und vollständiger Zielraum &
Gleichmäßiger Grenzwert, auch ohne Beschränktheit;
Satz~\ref{b47guide:thm:uniform-cauchy}.\\[4pt]
Strikte Abstandsverkleinerung auf nichtleerem kompaktem Raum &
Eindeutiger Fixpunkt und konvergente Iteration;
Satz~\ref{b47guide:thm:edelstein}.\\[4pt]
Beliebiger nichtleerer metrischer Raum &
Isometrische Darstellung durch Abstandsfunktionen;
Satz~\ref{b47guide:thm:kuratowski-frechet}.\\[4pt]
Vollständiger metrischer Raum &
Vollständiger Raum seiner nichtleeren kompakten Teilmengen;
Satz~\ref{b47guide:thm:hausdorff-complete}.\\[4pt]
Kompakter metrischer Raum &
Hausdorff-konvergente Teilfolgen kompakter Mengen;
Satz~\ref{b47guide:thm:blaschke}.\\
\bottomrule
\end{tabularx}
\end{center}
\section{Wichtige Unterschiede auf einen Blick}

\begin{center}
\small
\begin{tabularx}{\linewidth}{@{}p{.40\linewidth}X@{}}
\toprule
Voraussetzung & Was daraus allein nicht folgt\\
\midrule
Beschränktheit & Totale Beschränktheit oder Kompaktheit.\\[4pt]
Vollständigkeit und Beschränktheit & Kompaktheit.\\[4pt]
Totale Beschränktheit & Vollständigkeit.\\[4pt]
Abgeschlossenheit im Umraum & Vollständigkeit bei unvollständigem Umraum.\\[4pt]
Gleiche offene Mengen zweier Metriken & Gleiche Cauchy-Folgen oder
Vollständigkeit.\\[4pt]
Stetigkeit einer Abbildung & Gleichmäßige Stetigkeit oder Erhaltung
von Cauchy-Folgen.\\[4pt]
Gleichmäßige Stetigkeit & Vollständigkeit des Bildes eines vollständigen
Raums.\\[4pt]
Punktweise Beschränktheit einer Funktionenfamilie & Punktweise relative
Kompaktheit bei allgemeinen metrischen Zielräumen.\\[4pt]
Strikte Abstandsverkleinerung bei verschiedenen Punkten & Ein globaler
Kontraktionsfaktor $q<1$ oder ein Fixpunkt in jedem vollständigen Raum.\\
\bottomrule
\end{tabularx}
\end{center}

\renewcommand{\bibname}{Quellen und Bandzuordnung}
\begin{thebibliography}{99}
\addcontentsline{toc}{chapter}{Quellen und Bandzuordnung}
\footnotesize
\setlength{\itemsep}{6pt}
\setlength{\parskip}{0pt}
\interlinepenalty=10000

\bibitem{project}
Martin Kunze: \emph{Die Grundlagen der Mathematik}, Band~47:
\emph{Metrische Räume und Vollständigkeit}.
Metrik, Kugeln, Teilräume, Konvergenz, Cauchy-Folgen, Vollständigkeit,
Stetigkeit, Isometrien und Folgenkompaktheit; ergänzend
\path{tex/b47/compactness/main-result.tex} für totale Beschränktheit
und die Kompaktheitsäquivalenzen.
Die Gestaltung folgt den Leitfäden zu den Bänden~3, 4, 5, 28 und~37.
Die übrigen Einträge ergänzen klassische metrische Standardresultate.

\bibitem{metric-notes}
Jiří Lebl: \emph{Basic Analysis: Introduction to Real Analysis},
Kapitel~7: \emph{Metric Spaces}, insbesondere Abschnitte~7.1--7.5.
Autorenfassung: metrische Grundbegriffe, Folgen, Vollständigkeit,
Kompaktheit und stetige Abbildungen.
\url{https://www.jirka.org/ra/html/ms_chapter.html}.

\bibitem{completion-notes}
Helga Baum: \emph{Grundkurs Analysis. Skript zur Vorlesung Analysis I--III},
Fassung vom 5.~Oktober 2012, Humboldt-Universität zu Berlin.
Kapitel~2: \emph{Metrische Räume}, insbesondere Definition~2.24 und
Satz~2.29 zur Vervollständigung; gedruckte Seite~62.
\url{https://www2.mathematik.hu-berlin.de/~wendl/Sommer2019/Analysis2/Baum_Analysis-BA-WS11-Summe.pdf}.

\bibitem{baire-notes}
Elmar Schrohe: \emph{Funktionalanalysis}, Kapitel~1:
\emph{Topologische und metrische Räume, Fortsetzungssatz, Satz von Baire},
Leibniz Universität Hannover.
Satz~1.22 (Kompaktheitscharakterisierung), Satz~1.26 (gleichmäßig stetige
Fortsetzung), Satz~1.28 (Cantorscher Durchschnittssatz) und Satz~1.29
(Baire). Hier werden nur die metrischen Fassungen verwendet.
\url{https://www.analysis.uni-hannover.de/fileadmin/analysis/Schrohe/Schrohe_Lehre/Funktionalanalysis/fa1.pdf}.

\bibitem{banach-notes}
Manfred Einsiedler und Andreas Wieser:
\emph{Analysis I und II. Vorlesungsskript 2021/2022},
ETH Zürich, Fassung vom 5.~März 2022.
Kapitel~10: \emph{Metrische Räume}, Abschnitt~10.3.1,
Satz~10.41 (Banachscher Fixpunktsatz).
\url{https://metaphor.ethz.ch/x/2023/fs/401-1262-07L/sc/Analysis-Skript.pdf}.

\bibitem{ascoli-notes}
Benjamin McKay: \emph{Notes on Metric Spaces},
University College Cork, 27.~Januar 2021.
Kapitel~2, Satz~2.3: \emph{Arzelá--Ascoli I}, gedruckte Seiten~14--15.
Die hier verwendete Fassung nimmt einen vollständigen metrischen
Zielraum an. Dauerhafter institutioneller Datensatz:
\url{https://hdl.handle.net/10468/13437}.

\bibitem{koerner}
T. W. Körner: \emph{A Stroll Through Completeness and Compactness},
Autorenfassung vom 27.~April 2026, University of Cambridge.
Abschnitte~1.2, 1.3, 1.5--1.7 und~4.3:
Vollständigkeit, Kontraktionen, abgeschlossene Mengen, Baire,
Vervollständigung und Kompaktheit.
\url{https://www.dpmms.cam.ac.uk/~twk10/Main.pdf}.

\bibitem{countable-notes}
University of California, Riverside:
\emph{Mathematics 205A, Introduction to Topology---I, Course Notes---Part One},
revidierte Fassung, Herbst 2014.
Abschnitt~VI.1, Satz~VI.1.1 und Propositionen~VI.1.4, VI.1.6,
S.~91--93: abzählbare Basis, Separabilität und abzählbare Teilüberdeckungen
in metrischen Räumen sowie Vererbung auf Teilräume.
\url{https://math.ucr.edu/~res/math145B-2015/gentop-notes.pdf}.

\bibitem{product-notes}
J. B. Cooper: \emph{Elementary Topology},
Johannes Kepler Universität Linz, Abschnitt~3.3:
\emph{Products of metric spaces}, S.~36--37.
Gewichtete Produktmetrik, koordinatenweise Konvergenz,
Cauchy-Eigenschaft und Vollständigkeit.
\url{https://www.jku.at/fileadmin/gruppen/194/Lehre/Skripta_Jim/topology.pdf}.

\bibitem{lipschitz-extension}
Urs Lang: \emph{Notes on Rectifiability}, ETH Zürich,
Version vom 15.~Dezember 2023.
Proposition~3.1, Fall $n=1$: McShane--Whitney-Fortsetzung mit derselben
Lipschitz-Konstante. Die obere und die untere Formel entstehen aus den
Fassungen für $f$ und $-f$.
\url{https://people.math.ethz.ch/~lang/rect_notes.pdf}.

\bibitem{edelstein}
M. Edelstein: \emph{On Fixed and Periodic Points Under Contractive Mappings},
Journal of the London Mathematical Society 37 (1962), 74--79.
\url{https://doi.org/10.1112/jlms/s1-37.1.74}.
Die hier verwendete Fassung einschließlich der Iterationskonvergenz ist
auch in Tomonari Suzuki, \emph{Fixed point theorems in metric spaces}
(japanischer Originaltitel), RIMS K\^oky\^uroku 2041 (2017), 31--37,
Satz~1, abgedruckt:
\url{https://www.kurims.kyoto-u.ac.jp/~kyodo/kokyuroku/contents/pdf/2041-05.pdf}.

\bibitem{uniform-cauchy}
Marek Biskup: \emph{MATH 131BH notes}, UCLA, Kapitel~45:
\emph{Uniform convergence}, Fassung vom 18.~Mai 2026.
Definition~45.6, Lemma~45.7 und Korollar~45.8, S.~267--268:
gleichmäßiges Cauchy-Kriterium für metrische Zielräume.
\url{https://www.math.ucla.edu/~biskup/131bh.1.26s/PDFs/ch45.pdf}.

\bibitem{kuratowski-notes}
Urs Lang: \emph{Introduction to Geometric Measure Theory},
ETH Zürich, 2005, revidiert am 26.~November 2018.
Lemma~4.1, S.~27: Kuratowski--Fréchet-Einbettung.
Hier wird der Zielraum ausschließlich als metrischer Funktionenraum
verwendet; eine Normraumtheorie wird nicht vorausgesetzt.
\url{https://people.math.ethz.ch/~lang/diab05_notes_rev.pdf}.

\bibitem{hausdorff-notes}
Anton Petrunin: \emph{Lectures on metric geometry},
Autorenfassung, Kapitel~4: \emph{Space of subsets}.
Definition~4.1, Beobachtung~4.2, Satz~4.8 sowie Lemma~4.11 und
Aufgabe~4.12: Hausdorffmetrik, Blaschkes Auswahlsatz und Vollständigkeit
des Raums nichtleerer kompakter Mengen.
\url{https://anton-petrunin.github.io/metric-geometry/tex/lectures.pdf}.

\bibitem{hausdorff-distance}
Eddie Aamari: \emph{Hausdorff Distance and Measurability},
Lehrmaterial 2022/2023, Proposition~2.1, S.~3:
Beschreibung des Hausdorffabstands durch Abstandsfunktionen.
Die Durchmesserabschätzung findet sich außerdem in Petrunin,
Aufgabe~4.3.
\url{https://www.math.ens.psl.eu/~eaamari/teaching/2022-2023/M2_Jussieu/Lesson%202.pdf}.
\end{thebibliography}

\begingroup
\footnotesize
\noindent\textbf{Hinweise zum Quellengebrauch.}
Die Aussagen sind in einheitlicher Notation zusammengefasst; sie sind keine
wörtlichen Übernahmen. Die Quellen dienen dem Begriffs- und Satzabgleich.
Die Randfälle leerer Räume werden durch die Konventionen dieses Leitfadens
festgelegt. Online-Quellen eingesehen am 24.~Oktober 2026.
Die Auswahl enthält keine neuen Forschungsergebnisse und ist keine
historische Gesamtliste aller metrischen Sätze.
\endgroup